% -----------------------------------------------------------------------------
% Conteudo AODV-EN -- TCC de Huakson Lima e Diogo dos Reis Almeida (IFG Campus Inhumas).
% -----------------------------------------------------------------------------
\chapter{CONCLUSÃO}\label{cap:conclusao}

Este trabalho propôs, implementou e avaliou o AODV-EN, uma adaptação do protocolo
AODV (RFC 3561) para redes mesh multi-hop construídas sobre o protocolo ESP-NOW em
módulos ESP32. O núcleo reativo do protocolo --- descoberta de rotas sob demanda,
tabelas de rota com números de sequência, lista de precursores, fila de dados
pendentes, confirmação fim-a-fim e invalidação de rotas por falha de enlace --- foi
implementado como componente reutilizável em linguagem C, independente do ESP-IDF,
e validado tanto em simulação quanto em hardware com dez nós ESP32 dispostos
em cadeia multi-hop (três a quatro saltos).

Quanto ao objetivo geral, o trabalho demonstrou a viabilidade de prover roteamento
multi-hop sobre o ESP-NOW por meio de uma adaptação fiel do AODV, contornando a
ausência de roteamento nativo, o limite de \emph{peers} e as restrições de difusão do
protocolo. Em relação aos objetivos específicos, foram analisadas as limitações
técnicas do ESP-NOW (Capítulo~\ref{cap:referencial}), projetadas e implementadas as
adaptações necessárias (Capítulo~\ref{cap:implementacao}), e avaliado o desempenho do
algoritmo frente a um algoritmo de referência baseado em \emph{flooding}
(Capítulo~\ref{cap:resultados}).

A avaliação experimental, conduzida no cenário C1 com dez nós em cadeia multi-hop (três
a quatro saltos) e trinta repetições por algoritmo, evidenciou a superioridade decisiva
do roteamento reativo no regime multi-hop: o AODV-EN entregou 94,1\% dos pacotes contra
apenas 13,9\% do \emph{flooding}, que colapsa por não dispor de rota nem de
retransmissão fim-a-fim sob enlaces reais a 2~dBm. O AODV-EN apresentou, ainda, cerca de
20\% menos transmissões por pacote entregue (54,2 contra 68,8) --- métrica justa de custo
de canal, pois a carga de controle normalizada clássica não contabiliza as
retransmissões por inundação do \emph{flooding}. Pela mesma lógica, embora o consumo
energético agregado bruto do AODV-EN seja maior, a energia estimada por pacote
efetivamente entregue é cerca de cinco vezes menor (0,69 contra 3,75~J/pacote): a inundação
consome energia em transmissões que, em sua maioria, não chegam ao destino. A análise de
escalabilidade por simulação reforça o quadro, indicando que a desvantagem de custo do
\emph{flooding} se consolida e cresce com o tamanho e o diâmetro da rede, além do limite
de alcance imposto pelo TTL fixo. O cenário de malha densa (C3) acrescentou a nuance mais
importante: com caminhos redundantes (dez nós, cinco saltos), o \emph{flooding} não colapsa
--- a redundância eleva sua entrega a 77,8\%, contra 13,9\% na cadeia fina ---, de modo que
a vantagem de entrega do AODV-EN, ainda significativa (91,4\% contra 77,8\%), se estreita,
embora o roteamento preserve menor ocupação de canal e menor energia por entrega. A
vantagem do roteamento sobre a inundação revela-se, assim, \emph{dependente da topologia}.
Conclui-se que o AODV-EN troca uma sobrecarga de controle aproximadamente constante pela
garantia de entrega multi-hop --- inviável por inundação pura em topologias esparsas ---,
tornando-se mais vantajoso à medida que a rede cresce e os caminhos se tornam escassos.

O cenário de falha (C4) complementou essa análise ao exercitar a auto-recuperação do
protocolo: sob a queda cíclica de um relé do caminho (15~s a cada 45~s), o AODV-EN
manteve 70,6\% de entrega, re-descobrindo a rota pelo relé alternativo após cada queda,
conforme evidenciado pela curva de entregas acumuladas. O resultado explicita o
\emph{trade-off} entre eficiência e resiliência: o roteamento reativo paga um custo de
detecção e re-descoberta na falha, ao passo que a inundação, sem rotas a manter, é
estruturalmente tolerante à perda de um nó, porém ineficiente em banda.

Um resultado metodológico relevante decorreu da revisão da estratégia de disseminação:
a avaliação preliminar de um transporte por unicast sequencial revelou um viés
introduzido pelo mecanismo de retransmissão automática do unicast ESP-NOW, ausente no
\emph{broadcast}. A adoção da difusão por \emph{broadcast} de enlace, tanto para o
RREQ do AODV-EN quanto para o algoritmo de referência, restabeleceu a simetria
metodológica da comparação e alinhou o \emph{baseline} à definição canônica de
\emph{flooding} da literatura.

Como limitações, destacam-se a cobertura de cenários --- a campanha em hardware cobre os
cenários C1 (cadeia de dez nós), C3 (malha densa de dez nós) e C4 (falha induzida, seis
nós), todos com trinta repetições, ficando o cenário C2 (árvore hierárquica) como trabalho
futuro sobre a mesma infraestrutura ---, o caráter estimado do consumo energético e a coleta sequencial das
campanhas do AODV-EN e do \emph{flooding}, com nova distribuição física dos nós entre
elas. A seleção de rota empregou a métrica híbrida (contagem de saltos combinada
com penalização por RSSI), ativa no firmware da campanha; a política de gerenciamento
de \emph{peers} por substituição LRU também está implementada
(Capítulo~\ref{cap:implementacao}). A quantificação isolada do ganho dessas extensões
frente à contagem de saltos pura, contudo, não foi objeto de experimento dedicado e
constitui trabalho futuro.
A segurança do protocolo, por fim, não foi exercitada: o AODV-EN opera, como o AODV
original, sob a premissa de nós honestos e sem autenticação das mensagens de controle
(Seção~\ref{seguranca-roteamento-ad-hoc}), de modo que o endurecimento contra os ataques
conhecidos e a avaliação de seu custo em ESP32 também ficam como trabalho futuro.

Como desdobramentos, propõem-se: a ampliação do número de repetições e a cobertura dos
demais cenários experimentais; a separação física dos nós para exercitar múltiplos
saltos reais; a medição física de energia com instrumentação dedicada; e a
avaliação experimental da política LRU de \emph{peers} e da métrica híbrida de
roteamento em cenários multi-hop reais. Tais extensões consolidariam o AODV-EN como uma alternativa padronizada,
de baixo custo e energeticamente eficiente por entrega útil para redes mesh sobre ESP-NOW.
